\subsection{CP-03: motivo de rechazo normalizado en Kushki y Mercado Pago}
\label{sec:cp03}
\label{anexo:cp03}

\FichaCaso
  {CP-03}
  {Equipo Kit Pagos Colombia}
  {09/11/2026}
  {Motivo de rechazo normalizado en Kushki y Mercado Pago}
  {Pendiente (\#174)}
  {Pendiente (\#174)}
  {Provocar rechazos con datos de prueba de Kushki y Mercado Pago y comprobar que el SDK devuelve el código nativo y la categoría normalizada en \texttt{rejectionReason}, que una transacción aprobada no tiene ese campo y que el motivo no se pierde al pasar por la ruta REST \texttt{/v1/api}. El caso cubre el defecto que producía un objeto \texttt{\{\}} cuando el esquema de serialización descartaba los campos.}
  {Funcional (sistema, caja negra sobre la ruta real).}
  {\textbullet\ Dependencias de \texttt{simulator-api/} instaladas, incluido el paquete \texttt{kit-pagos-colombia} que usa la suite.\newline
   \textbullet\ No se necesita \texttt{.env}: la suite levanta la app en memoria y usa credenciales ficticias.\newline
   \textbullet\ La suite no abre un puerto externo; usa \texttt{app.inject()}.}
  {Local, en memoria. Node.js 20; API de Simulación construida con \texttt{buildApp()}; sin red externa ni credenciales reales. Los datos de la tarjeta y del titular son ficticios.}

\subsubsection{Objetivo}

Comprobar que el motivo de un rechazo llega completo al cliente con dos campos: \texttt{rejectionCode} (el código de la pasarela) y \texttt{rejectionCategory} (la categoría normalizada del SDK). El caso falla si el campo llega ausente, vacío o como \texttt{\{\}}, y también si una transacción aprobada lo incluye.

\subsubsection{Procedimiento}

\begin{TablaPasos}{Pasos del caso CP-03}
  1 &
  Ejecutar la suite de motivo de rechazo. &
  \texttt{cd simulator-api \&\& npx jest rejection-reason --coverage=false}. &
  La suite reporta 5 pruebas en 1 archivo, todas en verde, y termina con código 0. &
  & & \\ \hline
  2 &
  Verificar el rechazo de Kushki por fondos insuficientes. &
  Tarjeta \texttt{4349001210846432}, monto \texttt{50000.00 COP}, tokenización por \url{/v1/sim/kushki/card/v1/tokens}. &
  \texttt{POST}\ \url{/v1/api/payments} responde 201; la transacción tiene estado \texttt{DECLINED} y \texttt{rejectionReason} igual a \url{{ rejectionCode: "021", rejectionCategory: "INSUFFICIENT_FUNDS" }}. &
  & & \\ \hline
  3 &
  Verificar un segundo código de Kushki. &
  Tarjeta \texttt{4349008516656431}. &
  Estado \texttt{DECLINED} y \texttt{rejectionReason} igual a \url{{ rejectionCode: "019", rejectionCategory: "UNKNOWN" }}. El código nativo se conserva aunque la categoría sea \texttt{UNKNOWN}. &
  & & \\ \hline
  4 &
  Verificar el rechazo de Mercado Pago por fondos insuficientes. &
  Titular \texttt{FUND Ruiz}; token de tarjeta creado en \url{/v1/sim/mercadopago/card_tokens}. &
  Estado \texttt{DECLINED} y \texttt{rejectionReason} igual a \url{{ rejectionCode: "insufficient_amount", rejectionCategory: "INSUFFICIENT_FUNDS" }}. &
  & & \\ \hline
  5 &
  Verificar el rechazo de Mercado Pago por emisor. &
  Titular \texttt{OTHE Ruiz}. &
  Estado \texttt{DECLINED} y \texttt{rejectionReason} igual a \url{{ rejectionCode: "rejected_by_issuer", rejectionCategory: "UNKNOWN" }}. &
  & & \\ \hline
  6 &
  Verificar que la consulta devuelve el mismo motivo. &
  El identificador de una transacción rechazada, con \texttt{GET}\ \url{/v1/api/payments/:id?gateway=mercadopago}. &
  Respuesta 200 con el mismo \texttt{rejectionReason} de la creación; el motivo no se pierde entre la creación y la consulta. &
  & & \\ \hline
  7 &
  Verificar que una transacción aprobada no tiene motivo. &
  Titular \texttt{APRO Ruiz}, tarjeta de Mercado Pago. &
  Estado \texttt{APPROVED} y la transacción no tiene la propiedad \texttt{rejectionReason}. &
  & & \\ \hline
\end{xltabular}
\end{TablaPasos}

\subsubsection{Criterio de aceptación}

\begin{enumerate}
  \item Las 5 pruebas de la suite pasan y el proceso termina con código 0.
  \item Los cuatro rechazos devuelven exactamente el par \texttt{rejectionCode}/\texttt{rejectionCategory} esperado, no un objeto vacío ni un campo ausente.
  \item La consulta por identificador conserva el motivo de la creación.
  \item Una transacción aprobada no expone \texttt{rejectionReason}.
  \item La verificación pasa por la ruta HTTP real y por el SDK instalado, no solo por la función de serialización.
\end{enumerate}

\subsubsection{Observaciones de diseño}

La suite existe por un defecto concreto: \texttt{fast-json-stringify} serializa la respuesta con el esquema de la ruta y descarta lo que el esquema no declara. Con \texttt{code}, \texttt{category} y \texttt{message} en el esquema y \texttt{rejectionCode} y \texttt{rejectionCategory} en el objeto del SDK, el motivo llegaba como \texttt{\{\}} a cualquier cliente. Por eso la prueba atraviesa la ruta real y el SDK instalado: una prueba de la función de serialización aislada no habría reproducido el problema.

El caso también comprueba que la categoría \texttt{UNKNOWN} es válida: el SDK no inventa una categoría cuando el código nativo no tiene un equivalente conocido, y conserva el código para que el comercio pueda mostrarlo. Es preferible un \texttt{UNKNOWN} con código que un motivo inventado.

\subsubsection{Trazabilidad}

RF-02, RF-05, RF-08, RF-11 y RF-18; HU-05, HU-07 y HU-11.
